\chapter{Four harder tests}\label{ch:deeper}

The first notebook ends with what it did not do. The second notebook, \texttt{climate\_data\_deeper.ipynb}, does those
four things: it goes back to about 27,000 weather stations, raw and corrected, and rebuilds the averages, including
from the study's own stations; it removes El Ni\~no and La Ni\~na; it runs every trend on all 200 versions of the global
record; and it compares the observed trends with full climate models read straight from the cloud. It needs, besides
the four libraries of the first notebook, \texttt{xarray}, \texttt{zarr} and \texttt{fsspec} (all in Colab), and
downloads about 230 MB the first time.

\section{The plan}
\input{generated/nb_climate_data_deeper_1}
\section{Test 1: back to the stations}
\input{generated/nb_climate_data_deeper_2}
\section{Test 2: El Ni\~no}
\input{generated/nb_climate_data_deeper_3}
\section{Test 3: the uncertainty of the record}
\input{generated/nb_climate_data_deeper_4}
\section{Test 4: what climate models say}
\input{generated/nb_climate_data_deeper_5}
\section{What the four tests show}
\input{generated/nb_climate_data_deeper_6}
\section{Recording what the run used}
\input{generated/nb_climate_data_deeper_7}

\section{Where to go next}
\begin{itemize}
\item Change the gridding from 5\textdegree{} to 2\textdegree{} or 10\textdegree{} boxes: how much do the trends move?
\item Use the daily archive for a handful of the study's stations, average maximum and minimum to monthly means, and compare with the monthly file.
\item Add more CMIP6 models, or several runs of one model, and see how much of the model spread is chance.
\end{itemize}
